\section{Introduction} \label{section:intro}
% importance of learning type system
Recently, language models (LMs) have achieved great success in code generation~\cite{li2022competition, DBLP:conf/pldi/Xu0NH22, zheng2023codegeex, DBLP:journals/corr/abs-2108-07732}, where code is generated based on a natural language prompt.
%
LMs typically treat programs as plain text, encoding them as sequences of tokens and generating code by iteratively predicting the next token given the current partial program. After being trained on large corpora of code and natural language, such LMs can generate programs on demand with high accuracy.

However, language models (LMs) frequently struggle with type systems, often producing code that violates basic type constraints. 
Empirical studies indicate that type errors alone account for 33.6\% of failed LM-generated programs~\cite{DBLP:journals/ese/TambonDNKDA25, DBLP:journals/corr/abs-2407-06153}. 
Notably, even cutting-edge tools powered by state-of-the-art LMs can generate code that fails to compile.
For instance, an evaluation of GitHub Copilot on LeetCode revealed that 24\% of its suggestions resulted in compilation errors~\cite{DBLP:conf/msr/NguyenN22}, which mainly resulted from type errors~\cite{DBLP:journals/corr/abs-2504-09246}. 
This type validity problem becomes particularly pronounced in resource-limited languages or those with complex type systems~\cite{DBLP:journals/corr/abs-2406-00515,DBLP:journals/corr/abs-2410-03981}.

%Theoretically, while enhancing type systems with additional symbolic constraints could improve code safety, such strengthened systems would inevitably suffer from even scarcer training data, creating a fundamental limitation for data-driven approaches. 
%Despite the challenge, recent research offers a promising benefit: by reducing the LMs' burden of learning these constraints implicitly, we can free up their capacity for mastering program semantics and algorithmic logic, thereby significantly boosting overall performance. This aligns with findings in syntactic learning, where explicit structural guidance has been shown to enhance correctness across model scales~\cite{zhang2025gramtransbettercoderepresentation, DBLP:conf/acl/LiangZ0LLXCZZZC25}.
Ensuring type correctness is also useful beyond reducing type errors. First, the type system is essentially a static analysis system, and in theory, we can incorporate any decidable constraint into the type system.
%It is better to give an important decidable property to astonish the reviewer.
For example, Rust ensured no invalid memory access through an ownership type system. Therefore, by enriching the type system, we can enforce any decidable safety property of the code. Second, as revealed by recent studies~\cite{zhang2025gramtransbettercoderepresentation, DBLP:conf/acl/LiangZ0LLXCZZZC25}, even if LMs can learn to satisfy simple code constraints such as syntactic correctness given sufficient training data, by assisting LMs through external facilities for satisfying the constraints during training and inference, the overall performance of the model will further improve, possibly because the computational resources of the model can then be saved for higher-level tasks, such as algorithm design and aligning the semantics between the programs and the natural language description. 

% challenges in avoiding type errors
Ensuring type correctness in practice, however, remains challenging. 
As illustrated in \autoref{fig:codet5-incorrect}, which presents an ill-typed output from a widely-used and fine-tuned model \toolname{CodeT5}~\cite{DBLP:conf/emnlp/0034WJH21}, the model incorrectly references an undefined variable \tcode{num}, which could have been prevented through proper type checking.
Essentially, the challenge originates from the \textit{huge gap} between the text-based token sequence representation of programs and the actual structural process of type checking.

\begin{figure}
\hfill
    \begin{subfigure}{0.45\textwidth}
        \begin{lstlisting}[language=java, basicstyle=\footnotesize\ttfamily, escapeinside={@}{@}, numbers=left, numbersep=2pt]
int firstOdd(List<Integer> l) {
  for (int i = 0; i < l.size(); i++)
    if (l.get(i) % 2 != 0)
      return @\greenbox{l.get(i)}@;
  return -1;
}
        \end{lstlisting}
\vspace{-0.5em}
\caption{The expected function.}
\vspace{1em}
        \label{fig:codet5-correct}
    \end{subfigure}
    \hfill
    \begin{subfigure}{0.45\textwidth}
\centering 
\begin{minipage}{0.6\textwidth}
\begin{lstlisting}[language=java, basicstyle=\footnotesize\ttfamily, escapeinside={@}{@}, numbers=left, numbersep=3pt, firstnumber=3]
if (l.get(i) % 2 != 0)
  return @\redbox{num}@;
\end{lstlisting}
\end{minipage}
\vspace{-0.5em}
\caption{The ill-typed code generated by \toolname{CodeT5}}
\label{fig:codet5-incorrect}
\vspace{1em}
\begin{minipage}{0.6\textwidth}
\begin{lstlisting}[language=java, basicstyle=\footnotesize\ttfamily, escapeinside={@}{@}, numbers=left, numbersep=3pt, firstnumber=3]
if (l.get(i) % 2 != 0)
  return @\redbox{i}@;
\end{lstlisting}
\end{minipage}
\vspace{-0.5em}
\caption{The still-incorrect code after incorporating constraint decoding for types.}
        \label{fig:codet5-constraint-incorrect}
    \end{subfigure}
\hfill\ 
\caption{A code generation task and the incorrect Java programs generated by \toolname{CodeT5}, where the prompt comprises the function signature, input-output examples, and the description ``\emph{return the first odd element from a list (-1 by default)}''. We highlight the incorrect code as \redbox{red} and the expected counterpart as \greenbox{green}.}
    \label{fig:intro-sample}
\end{figure}

\begin{itemize}
    \item During training, LMs can only see sequences of syntax tokens, which provides little information on the underlying type system. From these sequences alone, LMs struggle to recover and learn the entire type system.
    %
    For example, typing the function in \autoref{fig:codet5-correct} relies on the signature of the \tcode{get} method and the auto-unboxing rule in Java that allows \tcode{Integer} to be returned as \tcode{int}, but both rules are obscured within the text-based token sequence.
    
    \item During generation, LMs must ensure the type correctness of the produced code. However, type-checking the partially generated code is challenging, as it often requires complex analysis to gather necessary contextual information from the existing code. 
    For instance, to check variable availability, the LM must identify all prior definitions (such as \tcode{l} and \tcode{i} in \autoref{fig:codet5-correct}) and determine whether each variable remains in scope, which necessitates non-trivial global reasoning across the entire codebase. Existing LMs often struggle with such complex contextual reasoning, especially as the codebase grows in size~\cite{guo_2025,DBLP:conf/acl/BaiTZ0WLCX0D0L25}.
\end{itemize}

To enforce type correctness, a typical approach is \emph{rejecting invalid generations}, as realized by a series of {constrained decoding} approaches~\cite{DBLP:conf/sigsoft/0003X023,DBLP:journals/corr/abs-2504-09246}.
These approaches leverage external type checkers and reject untypable programs generated by LMs.
%
However, such an external patch cannot bridge the gap in principle.
%
Constrained decoding can only exclude ill-typed programs; it does not help the model learn type systems, and thus cannot improve, and may even worsen, the distribution among well-typed programs.~\cite{DBLP:conf/nips/ParkWBPD24}.
%
For example, in our sample task, \toolname{CodeT5} assigns a low probability to the expected expression (\autoref{fig:codet5-correct}), possibly because the LM is unaware of the auto-unboxing rule and thus incorrectly identifies it as ill-typed.
%
Constrained decoding helps little with this case, as shown in \autoref{fig:codet5-constraint-incorrect} -- it merely shifts to a still incorrect but well-typed program that avoids auto-unboxing.

To guide the model in learning the type system, a straightforward approach is to employ a reasoning model that augments the program with a type correctness proof; that is, the model first generates the program and then generates its type correctness proof as the type reasoning process, or vice versa. In this way, the type reasoning is made explicit to the model, enabling it to learn the type system of the target program.
%
However, the separation between program generation and type reasoning introduces fundamental limitations. During training, presenting them as separate phases prevents LMs from capturing their fine-grained interactions, requiring extra effort to align them. During inference, the model must still perform global type reasoning without localized typing context, leaving the challenge of contextual reasoning unresolved.
%
As a result, this approach not only hinders learning efficiency but also causes inconsistencies between type reasoning and generated code. As our evaluation will later demonstrate, it does not yield satisfactory performance.

% our approach
Building on these insights, we aim to propose a new synthesis system that \emph{reconciles type systems internally with LMs}. 
Our approach leverages the fact that LMs are not confined to generating programs as plain text; they can be incorporated into any search-based system that constructs programs via sequences of decisions, with each decision made in a specific context. 
For any such system, an LM can be trained to predict the optimal decision at each step, enabling program synthesis by interacting with the LM throughout the search process.

In this way, reconciling type systems with LMs reduces to designing a synthesis system $\mathcal S$ that cooperates effectively with LMs while ensuring type correctness.
Specifically, we expect the system $\mathcal S$ to satisfy the following properties, which succinctly summarize the main challenges of LM-based code generation as we discussed earlier:
\begin{itemize}
 \item \textbf{Type Explicitness}. $\mathcal{S}$ should explicitly trace the type derivation throughout the decision process, making the type reasoning visible to the model.
 \item \textbf{Context Locality}. At each step, $\mathcal{S}$ should present all necessary type information within a small, bounded context window. This relieves LMs from the burden of inferring such information from lengthy contexts.
 \item \textbf{Derivation Vicinality}. The decisions in $\mathcal{S}$ should be represented such that each program fragment appears adjacent to its type derivation, making it easier for LMs to learn the correspondence between code and type reasoning during training.
 \item \textbf{Data Usability}. The decision sequences of $\mathcal{S}$ and source code should be automatically and bidirectionally convertible, enabling both the extraction of decision sequences from existing programs and the reconstruction of programs from generated decision sequences.
\end{itemize}

\textbf{\textit{Our first contribution}} is the observation that, under constructive logic, an existential proof of type correctness inherently contains the program it proves. Formally, an existential proof $\Pi$ of the following statement
\[
\exists p.\; \mathsf{welltyped}(p)
\]
must explicitly construct a witness program $p$. Thus, there exists an extraction function $\mathsf{extract}$ such that $\mathsf{extract}(\Pi) = p$. We leverage this insight to design a novel synthesis system and program representation that unifies program generation with type reasoning.

Based on this observation, the process of constructing an existential proof can be reformulated as a program synthesis process. The proof is built by applying typing rules step by step, where each step makes a \textit{decision} that either applies a rule to decompose the current goal into subgoals, or instantiates a variable within the context. Since the program is embedded within the proof, this sequence of decisions simultaneously determines the program being synthesized.

We use this decision sequence as a novel program representation, which satisfies the desired properties:
\begin{itemize}
    \item \textbf{Type Explicitness}. Constructing the decision sequence is equivalent to constructing the existential type correctness proof, making the complete type reasoning explicit throughout the process.
    \item \textbf{Derivation Vicinality}. The representation inherently captures both the program and its type correctness proof simultaneously, as they arise from the same construction process.
\end{itemize}

Notably, all proof rules in our system are formulated as constrained Horn clauses (see \autoref{chc-intro}). This representation not only subsumes the prior work on syntactic correctness~\cite{DBLP:conf/icse/ZhuL000C24,DBLP:conf/acl/LiangZ0LLXCZZZC25} but also accommodates a broader class of Turing-computable specifications. While this paper focuses on type constraints, our synthesis method is general and applicable to arbitrary logical constraints.

\textbf{\textit{Our second contribution}} is a novel encoder-decoder model architecture tailored for our synthesis framework. During synthesis, the current proof state evolves at each step, and the model needs to extract the dynamic typing context from this state to make the next decision. Our architecture addresses this by using the encoder to process both the natural language specification and the current synthesis goal, while the decoder autoregressively generates the decision sequence. This design enables the model to access localized type information at each step without requiring global reasoning over the entire codebase, thereby satisfying the \textbf{Context Locality} property.

\textbf{\textit{Our third contribution}} is \mainname, an automated system for training LMs to generate well-typed programs. \mainname takes as input (1) the language definition with its syntax, typing rules, and associated toolchain (including a parser and a type checker), and (2) a set of training tasks comprising prompts and target programs. 
From these inputs, \mainname automatically instantiates the synthesis framework for the target language, providing: (1) a training component that extracts decision sequences from programs and their existential type correctness proofs, and (2) a generation component that first generates decision sequences within the synthesis system, and then extracts well-typed programs from them.
This satisfies the \textbf{Data Usability} property: any well-typed program permits the extraction of its corresponding decision sequence, while any valid decision sequence can be reconstructed into a well-typed program along with its proof.

\smallskip
Our experimental results show that, relative to vanilla code generation, \mainname not only eliminates untypable programs\footnote{with respect to the type system we implemented}, but also substantially increases the number of programs that pass all unit tests. These findings indicate that \mainname successfully helps the LM acquire a deeper understanding of the type system and yields a more precise distribution over correct programs.

The rest of the paper is organized as follows. 
\autoref{sec:overview} gives an overview of our approach. 
\autoref{sec:meta} introduces how we translate the typing rules into synthesis rules. 
\autoref{sec:approach} proves the correctness of our framework and introduces how we prepare the training data and perform synthesis within the framework. 
\autoref{sec:implementation} introduces our neural network architecture. 
\autoref{sec:evaluation} describes the evaluation of our approach. 
Finally, \autoref{sec:related} discusses related works and \autoref{sec:conclusion} concludes the paper. 